% 三相混合最优潮流
% 本章文件由 \input 引入，不含导言区。
\section{三相混合最优潮流}\label{sec:threephase}

本章介绍平台的三相混合最优潮流模块（命名空间 \srcpath{hacdcpf::opf::phase\_hybrid}）。
该模块面向含单相/两相接入的不平衡配电系统与交直流混合系统，在\textbf{相域}
（phase domain）内直接建立并求解最优潮流：每个母线的每一相都是独立的电气节点，
支路以 $3\times 3$ 耦合导纳块建模，换流器按相端口接入并可选择相域控制律。
模块由三层组成——装配层 adapter 把 rich 工程模型 \fld{HybridPowerSystem} 装配为
相域 OPF 问题；OPF 本体把该问题作为单体非线性规划（NLP）用内点法求解；松弛层
给出同一问题的 SOCP 多面体外逼近与目标下界证书。与第~\ref{sec:acopfnative}~章的
单相 AC OPF 不同，本章模块不经过 canonical 投影与 \fld{SolverData} 装配，而是使用
专用的 \fld{ThreePhaseHybridOPFCase} 数据结构。

\textbf{成熟度声明（诚实口径）}：该模块是平台当前的活跃研发对象，接口、控制律覆盖
与验证深度仍在演进。本章全部结论以当前源代码为准；代码中的近似、占位、被拒绝的
输入与尚未生效的参数在~\ref{subsec:tp-limit}~节逐条列出，不作为已实现能力陈述。

\subsection{相域建模动机与适用范围}

\subsubsection*{动机}
输电网 OPF 的正序（平衡）等值在配电系统中失效：单相与两相分支、分相负荷、
分相分布式电源以及不平衡运行工况要求按相分别表示电压、电流与功率。代码层面
的对应设计如下：

\begin{itemize}
  \item \textbf{按相存在的节点}。装配层只为“母线在运且该相存在”的
        $(\text{母线},\varphi)$ 对分配相节点，\fld{phase\_mask} 不含某相时该相
        根本不进入问题，因而单相/两相分支自然得到少于 $3\times n_{\mathrm{bus}}$
        的节点集（\srcpath{src/optimal\_power\_flow/three\_phase\_hybrid\_adapter.cpp:build\_three\_phase\_hybrid\_model}）。
        节点的相归属记录在 \fld{ac\_phase\_index}（:build\_three\_phase\_hybrid\_model）。
  \item \textbf{分相负荷与分相电源}。母线三相负荷 \fld{pd\_a\_mw}/\fld{pd\_b\_mw}/%
        \fld{pd\_c\_mw}（无功同理）按相累加到对应节点的恒功率注入
        （:build\_three\_phase\_hybrid\_model）；三相发电机可按显式分相调度值或按相数均分
        接入单相节点（:build\_three\_phase\_hybrid\_model）。
  \item \textbf{相间耦合}。三相支路不以序分量解耦，而是以满 $3\times 3$ 复导纳块
        直接进入相域导纳矩阵 \fld{y\_ac}（:build\_three\_phase\_hybrid\_model），因此相间互感、
        不平衡潮流分布无需任何序网假设。
  \item \textbf{不平衡度约束}。对三相齐全的母线显式施加负序/正序电压幅值比
        （VUF）上限（\srcpath{src/optimal\_power\_flow/three\_phase\_hybrid\_opf.cpp:evaluate\_inequalities}），
        这是序分量在模块中的唯一显式用途（另见换流器正序电流控制）。
  \item \textbf{换流器相域控制}。VSC 换流器在相域内可选择均分相功率、
        跟网型正序电流、构网型下垂三种控制方程
        （\srcpath{include/hacdcpf/optimal\_power\_flow/three\_phase\_hybrid\_opf.hpp:PhaseVSCControlMode}），
        支撑不平衡接入下的逆变器调度研究。
\end{itemize}

\subsubsection*{适用范围（按装配层校验）}
装配入口 \srcpath{build\_three\_phase\_hybrid\_model} 对输入系统做硬性校验
（\srcpath{src/optimal\_power\_flow/three\_phase\_hybrid\_adapter.cpp:build\_three\_phase\_hybrid\_model}）：
必须含 \fld{three\_phase\_ac} 子系统、至少一个 DC 母线与一台 VSC 换流器；
\textbf{不支持} DC/DC 变换器与能量路由器，遇到即抛 \fld{std::invalid\_argument}。
此外还要求：至少一个在运相节点（:build\_three\_phase\_hybrid\_model）、至少一个交流参考源
（:build\_three\_phase\_hybrid\_model）、至少一个在运 DC 母线（:build\_three\_phase\_hybrid\_model）、DC 支路电阻
严格为正（:build\_three\_phase\_hybrid\_model）、每台 VSC 的交流三端子相 A/B/C 齐全
（:build\_three\_phase\_hybrid\_model）、至少一台两端均在运的 VSC（:build\_three\_phase\_hybrid\_model）。
可控 DC 静态发电机同样被拒绝（:build\_three\_phase\_hybrid\_model）。

\subsection{三层结构与代码组织}\label{subsec:tp-layers}

图~\ref{fig:tp-layers}~给出模块的三层组织与数据流。三层共享同一份问题数据
\fld{ThreePhaseHybridOPFCase}：装配层负责生成它，OPF 本体与松弛层各自独立地
消费它，二者不互相调用。

\begin{figure}[H]
\centering
\begin{tikzpicture}[
  box/.style={draw,rectangle,minimum width=5.2cm,minimum height=1.05cm,align=center,font=\footnotesize},
  arr/.style={-{Stealth},thick}]
  \node[box] (rich) {\fld{HybridPowerSystem}\\（rich 工程语义层）};
  \node[box,below=0.85cm of rich] (adapter) {装配层 adapter\\\srcpath{three\_phase\_hybrid\_adapter.cpp}};
  \node[box,below=0.85cm of adapter] (casedata) {\fld{ThreePhaseHybridOPFCase}\\（相域 OPF 问题数据）};
  \node[box,below left=1.0cm and -1.6cm of casedata] (opf) {相域 OPF 本体（NLP）\\\srcpath{three\_phase\_hybrid\_opf.cpp}};
  \node[box,below right=1.0cm and -1.6cm of casedata] (relax) {SOCP 松弛层（LP 序列）\\\srcpath{three\_phase\_hybrid\_relaxation.cpp}};
  \node[box,below=0.85cm of opf] (nlp) {\fld{engine::NLPModel}\\NativeIPM / Ipopt 后端};
  \node[box,below=0.85cm of relax] (lp) {\fld{engine::LPModel}\\HiGHS 后端};
  \draw[arr] (rich) -- (adapter);
  \draw[arr] (adapter) -- (casedata);
  \draw[arr] (casedata.south) -- (opf.north);
  \draw[arr] (casedata.south) -- (relax.north);
  \draw[arr] (opf) -- (nlp);
  \draw[arr] (relax) -- (lp);
\end{tikzpicture}
\caption{三相混合 OPF 模块的三层组织：adapter 负责 rich 模型到相域问题数据的装配，
OPF 本体与松弛层共享同一份 \fld{ThreePhaseHybridOPFCase}，分别下沉到
MIPSolvers 的 NLP 与 LP 求解通道}
\label{fig:tp-layers}
\end{figure}

\subsubsection*{装配层（adapter）}
\compmeta{ThreePhaseHybridModel}{\srcpath{include/hacdcpf/optimal\_power\_flow/three\_phase\_hybrid\_adapter.hpp}}
入口函数 \srcpath{build\_three\_phase\_hybrid\_model}
（头文件声明：\srcpath{include/hacdcpf/optimal\_power\_flow/three\_phase\_hybrid\_adapter.hpp:build\_three\_phase\_hybrid\_model}；
实现：\srcpath{src/optimal\_power\_flow/three\_phase\_hybrid\_adapter.cpp:build\_three\_phase\_hybrid\_model}）。
输出结构 \fld{ThreePhaseHybridModel}（\srcpath{include/hacdcpf/optimal\_power\_flow/three\_phase\_hybrid\_adapter.hpp:ThreePhaseHybridModel}）内嵌装配完成的
\fld{ThreePhaseHybridOPFCase}，并保留稳定组件 ID 与求解位置的映射：
\fld{ac\_bus\_ids}、\fld{bus\_phase\_nodes}（每母线三相的节点号，缺相为 $-1$）、
\fld{dc\_bus\_ids}、\fld{vsc\_component\_indices}、\fld{generator\_bindings}
（含组件 \fld{index}、母线、相号、名称、来源类型与爬坡参数）以及
\fld{model\_limitations} 字符串列表。查询接口 \texttt{phase\_node(ac\_bus\_id, phase)}
按 ID 反查节点号（实现：src/optimal\_power\_flow/three\_phase\_hybrid\_adapter.cpp:phase\_node）。装配选项
\fld{ThreePhaseHybridAdapterOptions} 的全部字段见~\ref{subsec:tp-adapter}~节末的
参数表。

\subsubsection*{相域 OPF 本体（opf）}
\compmeta{ThreePhaseHybridOPFCase}{\srcpath{include/hacdcpf/optimal\_power\_flow/three\_phase\_hybrid\_opf.hpp}}
公开接口（\srcpath{include/hacdcpf/optimal\_power\_flow/three\_phase\_hybrid\_opf.hpp:solve\_three\_phase\_hybrid\_opf、make\_three\_phase\_hybrid\_pf\_case、solve\_three\_phase\_hybrid\_opf\_sequence、solve\_three\_phase\_hybrid\_opf\_branches}）：单点求解 \srcpath{solve\_three\_phase\_hybrid\_opf}
（实现入口 src/optimal\_power\_flow/three\_phase\_hybrid\_opf.cpp:solve\_three\_phase\_hybrid\_opf，核心实现 \srcpath{solve\_three\_phase\_hybrid\_opf\_impl}
于 src/optimal\_power\_flow/three\_phase\_hybrid\_opf.cpp:solve\_three\_phase\_hybrid\_opf\_impl）、运行点转潮流算例 \srcpath{make\_three\_phase\_hybrid\_pf\_case}
（src/optimal\_power\_flow/three\_phase\_hybrid\_opf.cpp:make\_three\_phase\_hybrid\_pf\_case）、固定拓扑序列求解 \srcpath{solve\_three\_phase\_hybrid\_opf\_sequence}
（src/optimal\_power\_flow/three\_phase\_hybrid\_opf.cpp:solve\_three\_phase\_hybrid\_opf\_sequence）与参数摄动分支求解 \srcpath{solve\_three\_phase\_hybrid\_opf\_branches}
（src/optimal\_power\_flow/three\_phase\_hybrid\_opf.cpp:solve\_three\_phase\_hybrid\_opf\_branches）。内部组织为：数据校验 \srcpath{validate\_case}
（src/optimal\_power\_flow/three\_phase\_hybrid\_opf.cpp:validate\_case）$\to$ 模型数据 \fld{ModelData} 与变量布局 \fld{Layout}
装配 \srcpath{build\_model\_data}（src/optimal\_power\_flow/three\_phase\_hybrid\_opf.cpp:build\_model\_data）$\to$ NLP 回调装配
\srcpath{build\_nlp}（src/optimal\_power\_flow/three\_phase\_hybrid\_opf.cpp:build\_nlp）$\to$ 初值/对偶初始化 $\to$ 内点求解
$\to$ KKT 审计与结果归因。

\subsubsection*{松弛/凸化层（relaxation）}
\compmeta{ThreePhaseHybridRelaxationResult}{\srcpath{include/hacdcpf/optimal\_power\_flow/three\_phase\_hybrid\_relaxation.hpp}}
入口 \srcpath{solve\_three\_phase\_hybrid\_opf\_relaxation}
（\srcpath{include/hacdcpf/optimal\_power\_flow/three\_phase\_hybrid\_relaxation.hpp:solve\_three\_phase\_hybrid\_opf\_relaxation}；实现：\srcpath{src/optimal\_power\_flow/three\_phase\_hybrid\_relaxation.cpp:solve\_three\_phase\_hybrid\_opf\_relaxation}），
以及同选项序列求解 \srcpath{solve\_three\_phase\_hybrid\_opf\_relaxation\_sequence}
（src/optimal\_power\_flow/three\_phase\_hybrid\_relaxation.cpp:solve\_three\_phase\_hybrid\_opf\_relaxation\_sequence）。头文件注释明确了该层性质：对相域 AC/DC SOCP 松弛做
多面体外逼近，每一条锥割都是 SOCP 可行集的必要条件，因此每一轮 LP 都保持为原
非凸 OPF 的外松弛；凸二次发电成本用支撑切线低估（include/hacdcpf/optimal\_power\_flow/three\_phase\_hybrid\_relaxation.hpp:solve\_three\_phase\_hybrid\_opf\_relaxation）。
该层输出的是\textbf{目标下界证书}而非可行运行点。

\subsection{装配层：rich 模型到相域问题数据}\label{subsec:tp-adapter}

\subsubsection*{标幺化与问题规模}
基准功率取 \fld{three\_phase\_ac.base\_mva}，缺失时回退 \fld{system.base\_mva}，
再缺失取 100~MVA（src/optimal\_power\_flow/three\_phase\_hybrid\_adapter.cpp:build\_three\_phase\_hybrid\_model）；VUF 上限 \fld{vuf\_max} 被截断到
$[0,1]$（:build\_three\_phase\_hybrid\_model）。相节点编号按母线顺序、相号 A/B/C 顺序稠密分配，
仅含在运母线的存在相（:build\_three\_phase\_hybrid\_model）。以下功率、电压除注明外均为该基准下
的标幺值。

\subsubsection*{三相导纳矩阵与节点数据}
\fld{y\_ac} 由配网三相潮流模块的 \srcpath{analysis::build\_full\_ybus\_phase\_entries}
生成（声明于 \srcpath{include/hacdcpf/power\_flow/distribution\_power\_flow.hpp}，
调用处 src/optimal\_power\_flow/three\_phase\_hybrid\_adapter.cpp:build\_three\_phase\_hybrid\_model；该函数细节超出本章范围）。调用前装配层复制一份
相域网络并清空其 \fld{external\_grids}——参考电压与可调度电源在 OPF 中已有
显式表示，若再压入其诺顿等值将重复计数（src/optimal\_power\_flow/three\_phase\_hybrid\_adapter.cpp:build\_three\_phase\_hybrid\_model）。
\fld{include\_shunts} 选项控制并联补偿是否入矩阵（:build\_three\_phase\_hybrid\_model）。
节点默认电压窗口为 $[0.90,1.10]$~pu（:build\_three\_phase\_hybrid\_model），随后被母线自身的
\fld{vmin\_pu}/\fld{vmax\_pu}（非正时回退默认）覆盖（:build\_three\_phase\_hybrid\_model）；
初值电压由母线 \fld{vm\_*\_pu} 与 \fld{va\_*\_deg}（度转弧度）合成
（:build\_three\_phase\_hybrid\_model）。三相齐全的母线登记进 \fld{three\_phase\_bus\_nodes}，
作为 VUF 约束的作用集（:build\_three\_phase\_hybrid\_model）。\fld{i\_ac\_fixed}（固定电流注入）
在装配层恒置零（:build\_three\_phase\_hybrid\_model）。

\subsubsection*{负荷等值（诚实口径）}
OpenDSS 风格负荷若 \fld{connection} 为 delta，或带有恒阻抗/恒电流成分
（\fld{const\_z\_percent} 等六个百分比字段任一非零），装配层\textbf{不能}直接
表示其电压特性，只能采用导入时的分相恒功率运行点等值；当选项
\fld{allow\_constant\_power\_load\_equivalent} 为 false 时这类负荷直接抛错，
为 true 时继续装配但在 \fld{model\_limitations} 中登记该近似
（src/optimal\_power\_flow/three\_phase\_hybrid\_adapter.cpp:build\_three\_phase\_hybrid\_model，限制说明字符串见:build\_three\_phase\_hybrid\_model）。

\subsubsection*{发电机与外部电网}
三相发电机逐相生成 \fld{PhaseGenerator}：若给出了显式分相调度
（$|P_a|+|P_b|+|P_c|+|Q_a|+|Q_b|+|Q_c|>10^{-12}$）则取分相值，否则把
\fld{p\_mw}/\fld{q\_mvar} 按相数均分；功率上下界同样按相数分摊，未给出有效
上下界（\fld{pmax\_mw}$\le$\fld{pmin\_mw}）时退化为固定出力（该相上下界相等，
src/optimal\_power\_flow/three\_phase\_hybrid\_adapter.cpp:build\_three\_phase\_hybrid\_model）。成本系数通过名称或母线在 rich 模型中匹配同机
（\srcpath{matching\_rich\_generator}，src/optimal\_power\_flow/three\_phase\_hybrid\_adapter.cpp:matching\_rich\_generator），匹配不到时
$c_2=0$、$c_1=30$（:build\_three\_phase\_hybrid\_model）；匹配到时二次项按相数放大
（\fld{cost\_c2}$\,{}^*=c_2\cdot n_\varphi$）。每台机组的每相随带一条
\fld{generator\_bindings} 记录用于结果归因（:build\_three\_phase\_hybrid\_model）。

外部电网（\fld{external\_grids}）产生两类对象：其一，按
\fld{use\_phase\_voltage\_setpoint} 选择分相或对称（$\theta$、$\theta-120^\circ$、
$\theta+120^\circ$）电压设定，登记为参考节点（src/optimal\_power\_flow/three\_phase\_hybrid\_adapter.cpp:add\_reference（lambda），对称角度
构造见:add\_reference（lambda））；其二，每相一台界为 $\pm 5\times$ 总负荷的虚拟发电机作为
外部电网功率交换的松弛变量（:build\_three\_phase\_hybrid\_model），成本同样来自 rich 侧匹配
（$c_2^*=3c_2$，缺省 $c_1=30$，:build\_three\_phase\_hybrid\_model）。若没有任何外部电网参考，
回退取首台 \fld{BusType::SLACK} 母线（:build\_three\_phase\_hybrid\_model）；仍无则抛错
（:build\_three\_phase\_hybrid\_model）。

\subsubsection*{直流网络与换流器端口}
DC 支路电导 $g=\max(1,\fld{n\_parallel})/\fld{r\_pu}$ 压入拉普拉斯型电导矩阵
\fld{g\_dc}（src/optimal\_power\_flow/three\_phase\_hybrid\_adapter.cpp:build\_three\_phase\_hybrid\_model，电导计算:build\_three\_phase\_hybrid\_model）。DC 母线自带
\fld{pd\_mw} 与电压初值/窗口（:build\_three\_phase\_hybrid\_model）；\fld{DC\_V} 型母线登记为
直流电压参考（:build\_three\_phase\_hybrid\_model）。DC 负荷按 \fld{p\_mw}$\times$\fld{scaling}
计入（:build\_three\_phase\_hybrid\_model）；不可控 DC 静态发电机、光伏阵列与 DC 储能以负负荷口径计入
（:build\_three\_phase\_hybrid\_model）；可控 DC 静态发电机被拒绝（:build\_three\_phase\_hybrid\_model）。若全系统未声明
\fld{DC\_V} 母线，第一台 VSC 的直流端子锚定直流电压，并登记
\fld{model\_limitations}（:build\_three\_phase\_hybrid\_model）。

每台在运 VSC 生成一个 \fld{PhaseVSC}：交流三端子按相号映射到相节点，缺任一相
即抛错（src/optimal\_power\_flow/three\_phase\_hybrid\_adapter.cpp:build\_three\_phase\_hybrid\_model）；效率截断到 $[0.01,1]$（:build\_three\_phase\_hybrid\_model）；容量
\fld{s\_max\_pu} 取 \fld{p\_rated\_mw} 与各功率限幅绝对值的最大者除以基准，
选项 \fld{enforce\_converter\_capacity} 关闭时取 $10^6$（:build\_three\_phase\_hybrid\_model）；相电流上限
取 \fld{i\_ac\_max\_pu}，选项 \fld{enforce\_converter\_current} 关闭或非正时取
$10^6$（:build\_three\_phase\_hybrid\_model）；\fld{ac\_grid\_forming} 或 \fld{AC\_GRID\_FORMING} 模式映射
为构网型下垂，其余映射为均分相功率（:build\_three\_phase\_hybrid\_model）；虚拟阻抗回退值
$r_v=0.01$、$x_v=0.10$~pu（:build\_three\_phase\_hybrid\_model）；电压参考取 \fld{v\_ac\_set\_pu}，非正回退
1.0（:build\_three\_phase\_hybrid\_model）。装配结束时登记总括性限制说明：adapter 仅支持相域交流导纳、
电阻型 DC 支路、直接 VSC 端口与恒功率节点负荷方程（:build\_three\_phase\_hybrid\_model）。

\begin{paramtable}
\fld{include\_shunts} & — & — & true/false & true & 并联补偿是否压入 \fld{y\_ac}（src/optimal\_power\_flow/three\_phase\_hybrid\_adapter.cpp:build\_three\_phase\_hybrid\_model）\\
\fld{enforce\_converter\_capacity} & — & — & true/false & true & 是否施加 VSC 视在功率上限，false 时 \fld{s\_max\_pu} 取 $10^6$（src/optimal\_power\_flow/three\_phase\_hybrid\_adapter.cpp:build\_three\_phase\_hybrid\_model）\\
\fld{enforce\_converter\_current} & — & — & true/false & true & 是否施加 VSC 相电流上限，false 或限值非正时取 $10^6$（src/optimal\_power\_flow/three\_phase\_hybrid\_adapter.cpp:build\_three\_phase\_hybrid\_model）\\
\fld{allow\_constant\_power\_load\_equivalent} & — & — & true/false & true & 允许 delta/电压相关负荷以恒功率运行点等值装配；false 时遇此类负荷抛错（src/optimal\_power\_flow/three\_phase\_hybrid\_adapter.cpp:build\_three\_phase\_hybrid\_model）\\
\fld{vuf\_max} & $\varepsilon_{\mathrm{vuf}}$ & pu & $[0,1]$（装配时截断） & 0.03 & 三相母线负序/正序电压幅值比上限（src/optimal\_power\_flow/three\_phase\_hybrid\_adapter.cpp:build\_three\_phase\_hybrid\_model）\\
\end{paramtable}

\subsection{相域数学模型}\label{subsec:tp-model}

\subsubsection*{决策变量与布局}
OPF 本体采用\textbf{直角坐标}电压表示：相节点 $k$ 的电压
$V_k=e_k+\mathrm{j}f_k$（\srcpath{src/optimal\_power\_flow/three\_phase\_hybrid\_opf.cpp:model\_voltage}，
函数 \srcpath{model\_voltage}）。原始决策向量按
\begin{align*}
x=\bigl[e;\,f;\,p_g;\,q_g;\,u;\,p^{ac};\,q^{ac};\,e^{gfm};\,f^{gfm};\,p^{dc}\bigr]
\end{align*}
排布：$e,f\in\mathbb{R}^{n_v}$ 为（降阶后）相节点电压实/虚部，$p_g,q_g\in\mathbb{R}^{n_g}$
为分相发电机出力，$u\in\mathbb{R}^{n_{dc}}$ 为直流节点电压，$p^{ac},q^{ac}\in\mathbb{R}^{n_{cp}}$
为换流器分相交流功率，$e^{gfm},f^{gfm}\in\mathbb{R}^{n_{gfm}}$ 为构网型换流器 A 相
内电势，$p^{dc}\in\mathbb{R}^{n_c}$ 为换流器直流功率（布局结构 \fld{Layout} 与各段
偏移：src/optimal\_power\_flow/three\_phase\_hybrid\_opf.cpp:Layout、482--499）。维度计数为
（\srcpath{src/optimal\_power\_flow/three\_phase\_hybrid\_opf.cpp:build\_model\_data}）：
\begin{align*}
n_{\mathrm{var}} &= 2n_v+2n_g+n_{dc}+2n_{cp}+2n_{gfm}+n_c,\\
n_{\mathrm{eq}} &= 2n_v+n_{dc}+n_c+n_{\mathrm{dcref}}+n_{\mathrm{ctrl}}+2n_{\mathrm{ref}},\\
n_{\mathrm{ineq}} &= 2n_{\mathrm{full}}+n_{3\varphi}+n_c+n_{cp},
\end{align*}
其中 $n_{\mathrm{ctrl}}$ 为换流器控制等式数：构网型每台 $2n_{cp,c}$ 条，其余
每台 $2(n_{cp,c}-1)$ 条（src/optimal\_power\_flow/three\_phase\_hybrid\_opf.cpp:build\_model\_data）；$n_{\mathrm{full}}$ 为降阶前相节点数，
$n_{3\varphi}$ 为三相齐全母线数。所有变量均为连续量，默认界 $\pm10^{20}$，
随后按物理意义收紧：$e,f\in[-1.2,1.2]$，$p_g,q_g$ 取机组界，
$u\in[\fld{v\_dc\_min\_pu},\fld{v\_dc\_max\_pu}]$，$p^{ac},q^{ac}\in[-s_{\max},s_{\max}]$
（定单位功率因数时 $q^{ac}=0$），$e^{gfm},f^{gfm}\in[-1.5,1.5]$，
$p^{dc}\in[-s_{\max},s_{\max}]$（\srcpath{src/optimal\_power\_flow/three\_phase\_hybrid\_opf.cpp:build\_nlp}，
函数 \srcpath{build\_nlp}）。

\subsubsection*{三相支路模型与节点功率平衡}
交流网络以 $n\times n$ 复稀疏耦合导纳矩阵 \fld{y\_ac} 建模（$n$ 为相节点数，
三相支路对应 $3\times 3$ 满块）。节点电流与复功率注入为
\begin{align*}
I &= Y_{ac}\,V + i_{\mathrm{fix}}, &
S_k &= V_k\,\overline{I_k},
\end{align*}
其中 $i_{\mathrm{fix}}$ 为固定电流注入向量 \fld{i\_ac\_fixed}（装配层恒零；
\srcpath{src/optimal\_power\_flow/three\_phase\_hybrid\_opf.cpp:evaluate\_equalities}，
函数 \srcpath{evaluate\_equalities}）。每个相节点的有功/无功平衡等式为
\begin{align*}
g_k &= \operatorname{Re}\{S_k\} + p^{\mathrm{load}}_k
  - \sum_{g\in\mathcal{G}(k)} p_{g} - \sum_{c\in\mathcal{C}(k)} p^{ac}_{c,\varphi(k)} = 0,\\
g_{n_v+k} &= \operatorname{Im}\{S_k\} + q^{\mathrm{load}}_k
  - \sum_{g\in\mathcal{G}(k)} q_{g} - \sum_{c\in\mathcal{C}(k)} q^{ac}_{c,\varphi(k)} = 0,
\end{align*}
$\mathcal{G}(k)$、$\mathcal{C}(k)$ 分别为挂在节点 $k$ 的发电机与换流器相端口
（\srcpath{src/optimal\_power\_flow/three\_phase\_hybrid\_opf.cpp:evaluate\_equalities}；
发电机项:evaluate\_equalities，换流器项:evaluate\_equalities）。解析雅可比按 $S_k$ 对 $(e,f)$ 的偏导
展开装配（\srcpath{equality\_jacobian}，src/optimal\_power\_flow/three\_phase\_hybrid\_opf.cpp:equality\_jacobian，其中导纳块偏导:equality\_jacobian）。

\subsubsection*{直流网络与换流器功率耦合}
直流侧以电导矩阵 \fld{g\_dc} 建模，节点 $k$ 的平衡等式与换流器耦合等式为
\begin{align*}
u_k\,(G_{dc}\,u)_k + p^{\mathrm{dc,load}}_k - \sum_{c\in\mathcal{C}_{dc}(k)} p^{dc}_c &= 0,\\
p^{dc}_c + \eta_c \sum_{\varphi} p^{ac}_{c,\varphi} &= 0,
\end{align*}
即直流注入等于电阻网络功率加换流器抽取，换流器直流功率与三相交流有功之和
以效率 $\eta_c$ 闭合（\srcpath{src/optimal\_power\_flow/three\_phase\_hybrid\_opf.cpp:evaluate\_equalities}；
直流参考电压锚定 $u_t=u_{\mathrm{ref},t}$ 见:evaluate\_equalities）。该耦合为单效率近似，
不区分运行点的损耗曲线。

\subsubsection*{换流器相域控制方程}
每台换流器按 \fld{control\_mode} 选择一组等式（装配位置 src/optimal\_power\_flow/three\_phase\_hybrid\_opf.cpp:evaluate\_equalities，
类型枚举 \fld{PhaseVSCControlMode} 见 \srcpath{include/hacdcpf/optimal\_power\_flow/three\_phase\_hybrid\_opf.hpp:PhaseVSCControlMode}）：

（i）\textbf{均分相功率}（\fld{EqualPhasePower}）：以首相为基准，
\begin{align*}
p^{ac}_{\varphi} - p^{ac}_{a} = 0,\qquad
q^{ac}_{\varphi} - q^{ac}_{a} = 0,\qquad \varphi\neq a
\end{align*}
（\srcpath{src/optimal\_power\_flow/three\_phase\_hybrid\_opf.cpp:evaluate\_equalities}）。

（ii）\textbf{跟网型正序电流}（\fld{GridFollowingPLL}，别名
\fld{PositiveSequenceCurrent}）：记 $a=e^{\mathrm{j}2\pi/3}$，对
$\varphi\in\{b,c\}$ 施加
\begin{align*}
s_{\varphi}\,V_a - r_{\varphi}\,V_{\varphi}\,s_a = 0,\qquad
r_b=a,\; r_c=a^{2},
\end{align*}
其中 $s_\varphi=p^{ac}_\varphi+\mathrm{j}q^{ac}_\varphi=V_\varphi\overline{I_\varphi}$。
在电压非零时该式等价于 $I_\varphi=\overline{r_\varphi}\,I_a$，即换流器只输出
正序电流（\srcpath{src/optimal\_power\_flow/three\_phase\_hybrid\_opf.cpp:evaluate\_equalities}；
雅可比:equality\_jacobian；拉格朗日黑塞:lagrangian\_hessian）。

（iii）\textbf{构网型下垂}（\fld{GridFormingDroop}）：引入 A 相内电势
$E_a=e^{gfm}+\mathrm{j}f^{gfm}$，三相内电势取正序组
$E_\varphi\in\{E_a,\,a^{2}E_a,\,aE_a\}$，对每相施加
\begin{align*}
\overline{E_\varphi}\,V_\varphi - |V_\varphi|^{2} - \overline{z_v}\,s_\varphi = 0,
\qquad z_v = r_v + \mathrm{j}x_v,
\end{align*}
在 $V_\varphi\neq 0$ 时等价于诺顿等值 $E_\varphi = V_\varphi + z_v I_\varphi$
（\srcpath{src/optimal\_power\_flow/three\_phase\_hybrid\_opf.cpp:evaluate\_equalities}；
雅可比:equality\_jacobian）。校验要求构网型换流器虚拟阻抗非零（\srcpath{validate\_case}，
src/optimal\_power\_flow/three\_phase\_hybrid\_opf.cpp:validate\_case），序感知控制（后两种）要求三端子恰为相 A/B/C 且
\fld{ac\_phase\_index} 完整（src/optimal\_power\_flow/three\_phase\_hybrid\_opf.cpp:validate\_case）。

\textbf{诚实口径}：\fld{PhaseVSC} 中的 \fld{p\_droop\_pu}、\fld{q\_droop\_pu}、
\fld{voltage\_reference\_pu}、\fld{nominal\_frequency\_hz} 等字段被复制进内部
映射（src/optimal\_power\_flow/three\_phase\_hybrid\_opf.cpp:build\_model\_data）但\textbf{不进入任何等式/不等式}——所谓“下垂”在 OPF
中仅体现为上述诺顿耦合，功率-频率/无功-电压下垂律与电压设定值尚未构成约束；
\fld{pll\_kp}/\fld{pll\_ki} 与 \fld{voltage\_integral\_gain} 仅用于解后动态
平衡点恢复（src/optimal\_power\_flow/three\_phase\_hybrid\_opf.cpp:recover\_dynamic\_equilibria）。

\subsubsection*{不平衡约束（VUF）}
记三相齐全母线的节点电压为 $(V_a,V_b,V_c)$，正、负序分量按
\begin{align*}
V_1 = \tfrac{1}{3}\bigl(V_a + aV_b + a^{2}V_c\bigr),\qquad
V_2 = \tfrac{1}{3}\bigl(V_a + a^{2}V_b + aV_c\bigr)
\end{align*}
定义（\srcpath{src/optimal\_power\_flow/three\_phase\_hybrid\_opf.cpp:sequence\_components}，
函数 \srcpath{sequence\_components}），并对每座三相母线施加
\begin{align*}
|V_2|^{2} \;\le\; \varepsilon_{\mathrm{vuf}}^{2}\,|V_1|^{2},
\end{align*}
$\varepsilon_{\mathrm{vuf}}$ 即 \fld{vuf\_max}（\srcpath{src/optimal\_power\_flow/three\_phase\_hybrid\_opf.cpp:evaluate\_inequalities}，
函数 \srcpath{evaluate\_inequalities}；雅可比:inequality\_jacobian）。该约束只施加于
\fld{three\_phase\_bus\_nodes} 中登记的齐全三相母线，两相/单相分支不参与。

\subsubsection*{不等式约束汇总}
全部非线性不等式以 $h(x)\le 0$ 形式给出（求值
\srcpath{evaluate\_inequalities}，src/optimal\_power\_flow/three\_phase\_hybrid\_opf.cpp:evaluate\_inequalities；雅可比
\srcpath{inequality\_jacobian}，src/optimal\_power\_flow/three\_phase\_hybrid\_opf.cpp:inequality\_jacobian；拉格朗日黑塞的参数化模板
\srcpath{build\_inequality\_hessian\_template}，src/optimal\_power\_flow/three\_phase\_hybrid\_opf.cpp:build\_inequality\_hessian\_template）：
\begin{align*}
\underline{v}_k^{2} - |V_k|^{2} &\le 0, &
|V_k|^{2} - \overline{v}_k^{2} &\le 0, &
&\forall\,k,\\
|V_2|^{2} - \varepsilon_{\mathrm{vuf}}^{2}|V_1|^{2} &\le 0, &
\sum_{\varphi}\bigl((p^{ac}_{c,\varphi})^{2}+(q^{ac}_{c,\varphi})^{2}\bigr) - s_{\max,c}^{2} &\le 0, &
&\forall\,c,\\
(p^{ac}_{c,\varphi})^{2}+(q^{ac}_{c,\varphi})^{2}
  - i_{\max,c}^{2}\,|V_{c,\varphi}|^{2} &\le 0, &
& &\forall\,c,\varphi,
\end{align*}
依次为相节点电压幅值平方窗口（:evaluate\_inequalities）、VUF（:evaluate\_inequalities）、换流器三相总
视在功率（:evaluate\_inequalities）、换流器分相电流（:evaluate\_inequalities；由 $p^2+q^2\le i_{\max}^2|V|^2$
给出，$|V|$ 为恢复后的全节点电压）。变量界见决策变量一节。换流器相电流上限
未配置时回退 $s_{\max}/\sqrt{n_{\varphi}}$（src/optimal\_power\_flow/three\_phase\_hybrid\_opf.cpp:build\_model\_data）。

\subsubsection*{目标函数}
目标为发电成本加小幅正则化（\srcpath{src/optimal\_power\_flow/three\_phase\_hybrid\_opf.cpp:objective}，
函数 \srcpath{objective}）：
\begin{align*}
f = \sum_{g}\Bigl[c_{2,g}\bigl(S_{\mathrm{base}}p_g\bigr)^{2}
  + c_{1,g}\bigl(S_{\mathrm{base}}p_g\bigr)\Bigr]
  + 10^{-4}\sum_{g} q_g^{2}
  + 10^{-3}\sum_{c,\varphi}\bigl((p^{ac}_{c,\varphi})^{2}+(q^{ac}_{c,\varphi})^{2}\bigr)
  + 10^{-4}\sum_{c} (p^{dc}_c)^{2},
\end{align*}
其中成本以有名值 MW 计（$p$ 为标幺，$S_{\mathrm{base}}=$ \fld{base\_mva}），
三项正则化权重 $10^{-4}/10^{-3}/10^{-4}$ 为源码硬编码、不可配置
（梯度 \srcpath{objective\_gradient}，src/optimal\_power\_flow/three\_phase\_hybrid\_opf.cpp:objective\_gradient；黑塞
\srcpath{objective\_hessian}，src/optimal\_power\_flow/three\_phase\_hybrid\_opf.cpp:objective\_hessian）。成本币种代码中未固定。

\subsubsection*{图降阶与电压恢复}
\fld{ModelVariant::GraphReduced} 下，所有无注入的纯传输相节点经稀疏 Kron
（Schur 补）降阶消去：带负荷/固定电流注入、发电机、换流器端子、参考节点
一律保留（\srcpath{src/optimal\_power\_flow/three\_phase\_hybrid\_opf.cpp:build\_model\_data}，
函数 \srcpath{build\_model\_data}；降阶核为
\srcpath{graph::reduce\_sparse\_kron}，受 \fld{reduction\_options} 控制）。
\fld{ModelVariant::Full} 下保留全部节点，恢复算子取单位阵（:identity\_recovery）。
全节点电压由恢复算子线性还原
\begin{align*}
V_{\mathrm{full}} = R\,V_{\mathrm{model}}
\end{align*}
（\srcpath{src/optimal\_power\_flow/three\_phase\_hybrid\_opf.cpp:full\_voltage}，
函数 \srcpath{full\_voltage}；恢复行缓存:build\_model\_data）。任一全节点的恢复行为空
表明该节点与参考源电气脱开，此时抛运行错误并提示先移除无源组件
（:build\_model\_data）。消去节点数记入结果 \fld{eliminated\_phase\_nodes}（:solve\_three\_phase\_hybrid\_opf\_impl）。
注意电压窗口、VUF 等不等式始终定义在\textbf{全节点}上（
\srcpath{evaluate\_inequalities} 先恢复全电压再求值，src/optimal\_power\_flow/three\_phase\_hybrid\_opf.cpp:evaluate\_inequalities），
因此降阶不削弱约束覆盖面。

\subsubsection*{初值构造}
初值由 \srcpath{build\_initial\_point} 生成（src/optimal\_power\_flow/three\_phase\_hybrid\_opf.cpp:build\_initial\_point）：电压取
\fld{voltage\_start}（:build\_initial\_point），换流器功率按直流侧初始平衡反推并均分
（:build\_initial\_point）；随后进行“交流潮流初始化 + 换流器控制投影”三轮交替
（:build\_initial\_point），最后按节点平衡回算发电机出力并截断到上下界（:build\_initial\_point）。
交流潮流初始化 \srcpath{initialize\_ac\_power\_flow}（src/optimal\_power\_flow/three\_phase\_hybrid\_opf.cpp:initialize\_ac\_power\_flow）先尝试
直接牛顿法（至多 20 次迭代，线搜索要求全节点电压落在 $(0.5,1.3)$~pu，
:initialize\_ac\_power\_flow）；失败则回退 $x$ 并改用 Z-bus 不动点延拓（20 个负荷步、每步
至多 80 次迭代，电压幅值下限保护 0.2~pu，:initialize\_ac\_power\_flow）。存在非参考节点上的
发电机时初始化直接跳过（返回 false，:initialize\_ac\_power\_flow；调用方忽略该返回值并继续，
:build\_initial\_point——属已知粗糙点）。换流器控制投影
\srcpath{project\_converter\_control}（src/optimal\_power\_flow/three\_phase\_hybrid\_opf.cpp:project\_converter\_control）按控制律把当前总功率
重分配到各相：构网型反解 $E_a$ 并按 $I_\varphi=(E_\varphi-V_\varphi)/z_v$
重算相功率（:project\_converter\_control），跟网型按 $I_a=\overline{S_{\mathrm{tot}}}/(3V_1)$
正序重分配（:project\_converter\_control）。

\subsection{求解流程与求解器接口}\label{subsec:tp-solve}

\subsubsection*{主流程}
\srcpath{solve\_three\_phase\_hybrid\_opf\_impl}（src/optimal\_power\_flow/three\_phase\_hybrid\_opf.cpp:solve\_three\_phase\_hybrid\_opf\_impl）的步骤为：
数据校验与模型装配 $\to$ \fld{engine::NLPModel} 回调装配 $\to$ 初值构造
$\to$（NativeIPM 路径 Phase I）有预算的原始可行性恢复
$\to$ 对偶/松弛初始化 $\to$（Phase II）MIPSolvers Native IPM
$\to$ 结果归因与 KKT 审计。这里的 Phase I/II 是“可行原--对偶状态构造/最优性
求解”的接口分层，不是平衡 AC OPF 的 $t=0\to1$ 目标同伦。求解器容差统一取自
\fld{options.tolerance}，acceptable 级别放宽 10 倍（src/optimal\_power\_flow/three\_phase\_hybrid\_opf.cpp:solve\_three\_phase\_hybrid\_opf\_impl）。

\subsubsection*{后端：NativeIPM 与 Ipopt}
\fld{SolverBackend::NativeIPM} 走 MIPSolvers 原生内点核：滤波器全局化、关闭
问题缩放，Phase II 的 primal/dual/complementarity 始终使用
\fld{options.tolerance} 最终容差；Phase I 只构造 central corridor。适配器为
\srcpath{engine::NativeIPMAdapter}，求解调用:solve\_three\_phase\_hybrid\_opf\_impl；核头文件引入
见 src/optimal\_power\_flow/three\_phase\_hybrid\_opf.cpp:solve\_three\_phase\_hybrid\_opf\_impl）。该路径先做原始可行性恢复
\srcpath{restore\_primal\_feasibility}。若模型声明了自由控制列 $F$，其补集为
state/basic 列 $B$，每轮只解稀疏方阵 Newton 校正
\begin{equation}
 J_B(x_k)\,\Delta x_B=-g(x_k),\qquad \Delta x_F=0,
 \label{eq:phase-one-basic-newton}
\end{equation}
并用 SparseLU；分区不可用或方阵分解失败时，对矩形 $J$ 直接做 SparseQR。
该设计依据 Nocedal--Wright（2006）第 11.1 节的约束 Newton 校正，避免稠密
物化，也避免正规方程 $JJ^T$ 的条件数平方与额外填充。Deuflhard（2011）
第 2.2--2.3 节的单调阻尼规则只接受原 pu 坐标最大违反量严格下降的步；不引入
障碍参数序列、filter、二阶校正、Hessian 解或弹性变量，因而不复制 Phase II
的早期内点迭代。

Phase I 同时受迭代、稀疏数值分解和回溯三个硬预算约束。墙钟预算在分解之间
检查；单次稀疏分解不可协作中断，故它是软截止时间，可能被最后一个不可分割
分解略微越过，而总分解次数仍给出有限工作上界。有限墙钟预算
$\fld{phase\_one\_time\_limit\_ms}>0$ 时绝不调用 Ipopt；旧
\fld{warm\_start\_with\_ipopt} 只在显式无限预算（该字段 $\le0$）时保留为兼容
路径。
默认 $\mu_0=0.1$、$\eta_a=1$、$\eta_p=0.1$、$\eta_c=0.5$。初始违反量高于
$\eta_a\mu_0$ 时以零分解退出 Phase I；否则恢复到
$\epsilon_p=\max(\mathrm{tol},\eta_p\mu_0)$。slack floor 使用剩余 primal
corridor 的一半，并令 $\nu_i=\mu_0/s_i$，同时控制扰动 primal residual、乘子
尺度和中心性。Phase I 在原始 pu 坐标下重新审计，并请求 MIPSolvers 独立的
\fld{central\_warm\_start} 契约；它不同于最终可行的
\fld{primal\_feasible\_start}。只有完整 primal/dual/slack、原坐标与扰动 primal
residual、中心性及乘子尺度全部通过才保留 $x_0$，否则自动回到普通冷初始化。
审计请求/接受状态和实际稀疏线性求解
后端分别由 \fld{phase\_two\_start\_requested}、
\fld{phase\_two\_start\_accepted} 与
\fld{phase\_two\_linear\_solver\_backend} 返回。
\fld{SolverBackend::Ipopt} 直接把同一 NLP 交给
\srcpath{engine::IpoptAdapter}（src/optimal\_power\_flow/three\_phase\_hybrid\_opf.cpp:solve\_three\_phase\_hybrid\_opf\_impl）。\textbf{门控说明}：Ipopt
后端依赖嵌入式 Ipopt（宏 \srcpath{HACDCPF\_HAVE\_IPOPT}）；未编译时适配器不抛错，
而是返回 \fld{Unavailable:} 前缀的失败状态（
\srcpath{MIPSolvers/src/engine/solver/external/adapters.cpp:IpoptAdapter::available}），结果即
不收敛。默认后端为 Ipopt（\srcpath{include/hacdcpf/optimal\_power\_flow/three\_phase\_hybrid\_opf.hpp:ThreePhaseHybridOPFOptions}）。

\subsubsection*{对偶初始化与拉格朗日黑塞}
NativeIPM 路径的对偶初值由 \srcpath{initialize\_primal\_dual\_start} 构造：
仅当初点不超过上述 $\epsilon_p$ corridor 时启用；
松弛变量由 $s=-h(x_0)$ 截断生成（:initialize\_primal\_dual\_start），不等式乘子按障碍参数
$\mu/s$ 恢复（:initialize\_primal\_dual\_start）；等式乘子优先利用
\fld{equality\_free\_columns} 划分做基解（仅当 $n_{\mathrm{var}}-n_{\mathrm{eq}}=n_c$
时以每台换流器首相 $q^{ac}$ 为自由列），否则对稳态方程做一次 SparseQR
最小二乘恢复。求解器未返回等式对偶时（如 Ipopt 路径），在同一解点上重跑一次
该初始化并仅当对偶残差不超过 $10^{-6}$ 时采纳（src/optimal\_power\_flow/three\_phase\_hybrid\_opf.cpp:solve\_three\_phase\_hybrid\_opf\_impl）。

Phase I 的可行点通常不是最优点，因此完整 stationarity
$r_{\mathrm{stat}}=\nabla f+J_g^T\lambda+J_h^T\nu$ 应保留控制空间 reduced
gradient 给 Phase II，而不应被 Phase I 人为消去。对声明了 $B/F$ 分区的模型，
dual 质量证书为
\begin{equation}
 r_{\mathrm{dual,fit}}=
 \frac{\|r_{\mathrm{stat},B}\|_\infty}
 {1+\max(\|\lambda\|_\infty,\|\nu\|_\infty)};
 \label{eq:phase-one-dual-fit}
\end{equation}
SparseQR 回退则报告对应最小二乘正规残差
$\|J_g r_{\mathrm{stat}}\|_\infty$ 的同尺度值。结果中的
\fld{phase\_one\_dual\_fit\_residual} 用此口径判定 warm start 质量，
\fld{initial\_dual\_residual} 继续诚实报告完整 stationarity，二者不可互换。
拉格朗日黑塞 \srcpath{lagrangian\_hessian}（src/optimal\_power\_flow/three\_phase\_hybrid\_opf.cpp:lagrangian\_hessian）按目标黑塞、
潮流等式双线性项（:lagrangian\_hessian）、换流器控制等式项（:lagrangian\_hessian）与参数化
不等式模板（:lagrangian\_hessian）装配。

在 360 个三相母线、1080 个相节点、2175 个变量的稀疏回归上，初始违反量
$10^{-3}$ 已在默认 $10^{-2}$ corridor 内，因此 primal Newton 为 0 次，只做
一次 sparse-basis dual fit；dual-fit 为 $8.33\times10^{-15}$，五次计时中位数
为 5.438~ms。相较旧的严格 Phase I（两次 Newton、三次总分解），这里消除了
两次不必要的 primal 分解而不放宽 Phase II 最终证书。

\subsubsection*{KKT 审计与收敛口径}
求解器报成功后，模块在解点上\textbf{重新}评估物理等式与\textbf{全部}不等式
（含被预言机省略者），并施加审计阈值：原始/对偶/互补残差、电压越限、换流器
越限均不超过 $10^{-6}$ 才把 \fld{converged} 置真（
\srcpath{src/optimal\_power\_flow/three\_phase\_hybrid\_opf.cpp:solve\_three\_phase\_hybrid\_opf\_impl}，
审计条件:solve\_three\_phase\_hybrid\_opf\_impl）。即求解器自身成功只是必要条件。
\fld{verify\_derivatives} 为真时另做中心差分（步长 $10^{-5}$）校核等式/不等式
雅可比与拉格朗日黑塞，误差上界写入结果（:solve\_three\_phase\_hybrid\_opf\_impl；乘子取固定的伪随机模式
:solve\_three\_phase\_hybrid\_opf\_impl）。

\subsubsection*{约束预言机}
\fld{use\_constraint\_oracle} 为真时启用精确约束预言机
（src/optimal\_power\_flow/three\_phase\_hybrid\_opf.cpp:solve\_three\_phase\_hybrid\_opf\_impl）：以初值处 $h_j(x_0)\ge -\max(m_{\mathrm{init}},m_{\mathrm{act}})$
的行、全部换流器约束行与 \fld{oracle\_seed\_rows} 构成初始强制集
（:solve\_three\_phase\_hybrid\_opf\_impl）；每轮求解受限问题后在全量不等式上筛查，把
$h_j(x)\ge -m_{\mathrm{act}}$ 的行补入（:solve\_three\_phase\_hybrid\_opf\_impl）；
\fld{oracle\_probe\_iterations} 为正值时先跑“预测轮”（限迭代数），活跃集
稳定后再跑精确校正轮（:solve\_three\_phase\_hybrid\_opf\_impl）；相邻轮之间传递原始点与对偶
（:solve\_three\_phase\_hybrid\_opf\_impl）。达到 \fld{oracle\_max\_rounds} 仍有新增行时判不收敛
（:solve\_three\_phase\_hybrid\_opf\_impl）。被省略行的最大违反量记入 \fld{max\_omitted\_inequality}
（:solve\_three\_phase\_hybrid\_opf\_impl），作为近似口径的公开声明。

\subsubsection*{序列与分支求解}
\srcpath{solve\_three\_phase\_hybrid\_opf\_sequence}（src/optimal\_power\_flow/three\_phase\_hybrid\_opf.cpp:solve\_three\_phase\_hybrid\_opf\_sequence）要求各
时段网络与设备拓扑完全一致（否则抛错，:solve\_three\_phase\_hybrid\_opf\_sequence），用各时段最大负荷构造一次
共享 \fld{ModelData}（:solve\_three\_phase\_hybrid\_opf\_sequence），逐时段传递上一步的原始点、等式/不等式对偶
与松弛作为热启动，并把上一步强制行集作为预言机种子（:solve\_three\_phase\_hybrid\_opf\_sequence）。
\srcpath{solve\_three\_phase\_hybrid\_opf\_branches}（src/optimal\_power\_flow/three\_phase\_hybrid\_opf.cpp:solve\_three\_phase\_hybrid\_opf\_branches）从已认证的
基准点出发做同样的参数摄动分支求解，基准未收敛则抛错（:solve\_three\_phase\_hybrid\_opf\_branches）。

\subsubsection*{与单相 AC OPF 的接口关系}
本章模块与第~\ref{sec:acopfnative}~章的单相 AC OPF Native IPM 的关系是
“同核不同模”：两者共享 MIPSolvers 的 \fld{engine::NLPModel} 抽象与原生内点核
（及嵌入式 Ipopt 适配器），但三相模块不复用单相的 \fld{SolverData} 装配、
总线类型（PQ/PV/松弛）机制与投影回传，而是自建相域单体 NLP。结果侧的接口有
两个：其一，\srcpath{make\_three\_phase\_hybrid\_pf\_case}（src/optimal\_power\_flow/three\_phase\_hybrid\_opf.cpp:make\_three\_phase\_hybrid\_pf\_case）把
收敛运行点翻译为三相混合潮流算例 \fld{powerflow::ThreePhaseHybridPFCase}——
换流器总功率固定、按控制模式与是否直流电压参考映射到
\fld{EqualPhasePQ}\allowbreak /\fld{EqualPhaseVdcQ}\allowbreak /\fld{GridFollowingPQ}\allowbreak /\fld{GridFollowingVdcQ}\allowbreak /\fld{GridFormingVoltage}\allowbreak /\fld{GridFormingVdc} 六类潮流
控制（:make\_three\_phase\_hybrid\_pf\_case），供独立潮流复核 OPF 运行点（要求已收敛且维度匹配，
:make\_three\_phase\_hybrid\_pf\_case）；其二，上文序列/分支接口支撑时序与 N-1 类参数扫描。结果归因回
rich 组件依赖 adapter 保留的 \fld{generator\_bindings} 与
\fld{vsc\_component\_indices}。

\subsection{SOCP 松弛与目标下界证书}\label{subsec:tp-relax}

\subsubsection*{提升变量}
松弛层把电压双线性项提升为新变量：交流侧对每个相节点引入
$W_{kk}=|V_k|^2$，对每对相关联节点引入
$W_{ij}=V_i\overline{V_j}=W^{R}_{ij}+\mathrm{j}W^{I}_{ij}$（$i<j$）；
直流侧引入 $X_{kk}=u_k^2$ 与 $X_{ij}=u_i u_j$。关联节点对取自 \fld{y\_ac} 非零
模式、三相齐全母线的相两两组合、PSD 团（clique）候选与参考节点两两组合
（\srcpath{src/optimal\_power\_flow/three\_phase\_hybrid\_relaxation.cpp:solve\_three\_phase\_hybrid\_opf\_relaxation}；
团候选由导纳矩阵列邻居与三相母线生成并剔除被包含者:solve\_three\_phase\_hybrid\_opf\_relaxation）。
变量界为 $W_{kk}\in[\underline{v}_k^2,\overline{v}_k^2]$（参考节点固定为
$|V_{\mathrm{ref}}|^2$）、$W^{R/I}_{ij}\in[-\overline{v}_i\overline{v}_j,\,
\overline{v}_i\overline{v}_j]$（双参考对固定为 $V_i\overline{V_j}$ 的对应值）、
$X_{kk}\in[\underline{u}_k^2,\overline{u}_k^2]$、$X_{ij}\in[\underline{u}_i
\underline{u}_j,\overline{u}_i\overline{u}_j]$（:solve\_three\_phase\_hybrid\_opf\_relaxation）。发电机变量
$p_g,q_g$ 按原界引入，$c_2>10^{-14}$ 者附加成本上图变量 $t_g$（:solve\_three\_phase\_hybrid\_opf\_relaxation）；
换流器变量 $p^{ac},q^{ac}\in[-s_{\max},s_{\max}]$（定单位功率因数时
$q^{ac}=0$）与 $p^{dc}$（:solve\_three\_phase\_hybrid\_opf\_relaxation）。

\subsubsection*{线性化后的平衡方程与不等式}
提升后节点功率平衡成为线性等式
\begin{align*}
\sum_{j}\operatorname{Re}\{\overline{Y_{kj}}W_{kj}\} + p^{\mathrm{load}}_k
  - \sum_{g\in\mathcal{G}(k)} p_g - \sum_{c\in\mathcal{C}(k)} p^{ac}_{c,\varphi(k)} &= 0,\\
\sum_{j}\operatorname{Im}\{\overline{Y_{kj}}W_{kj}\} + q^{\mathrm{load}}_k
  - \sum_{g\in\mathcal{G}(k)} q_g - \sum_{c\in\mathcal{C}(k)} q^{ac}_{c,\varphi(k)} &= 0,\\
\sum_{j} (G_{dc})_{kj}X_{kj} + p^{\mathrm{dc,load}}_k
  - \sum_{c\in\mathcal{C}_{dc}(k)} p^{dc}_c &= 0,
\end{align*}
装配时按 $i<j$ 定向修正 $W^{I}$ 项符号（
\srcpath{src/optimal\_power\_flow/three\_phase\_hybrid\_relaxation.cpp:solve\_three\_phase\_hybrid\_opf\_relaxation}；
定向处理:solve\_three\_phase\_hybrid\_opf\_relaxation）。换流器耦合 $p^{dc}_c+\eta_c\sum_\varphi p^{ac}_{c,\varphi}=0$
与均分相功率等式按原样保留（:solve\_three\_phase\_hybrid\_opf\_relaxation）。VUF 约束对 $W=VV^{H}$ 是线性的：
$|V_2|^2-\varepsilon_{\mathrm{vuf}}^2|V_1|^2=\mathrm{tr}\bigl((n_2n_2^{H}-
\varepsilon_{\mathrm{vuf}}^2 n_1n_1^{H})W\bigr)\le 0$，其中 $n_1,n_2$ 为序分量
系数向量（:solve\_three\_phase\_hybrid\_opf\_relaxation）。若 \fld{y\_ac}（厄米特部最小特征值不小于
$-10^{-8}\cdot\max(1,\lVert Y\rVert)$）或 \fld{g\_dc} 通过无源性检查
（:ac\_admittance\_is\_passive），各加一条聚合网损非负割：发电与换流器注入之和不小于总负荷
（:solve\_three\_phase\_hybrid\_opf\_relaxation；未通过检查时放弃该割并登记 \fld{model\_limitations}）。

\subsubsection*{锥割与外逼近循环}
精确条件 $|W_{ij}|^2\le W_{ii}W_{jj}$（交流）、$X_{ij}^2\le X_{ii}X_{jj}$
（直流）、$\sum_\varphi(p^2+q^2)\le s_{\max}^2$ 与
$p_\varphi^2+q_\varphi^2\le i_{\max}^2 W_{nn}$ 均为旋转二阶锥约束；松弛层不调用
锥求解器，而是在每个 LP 解点上对违例锥生成支撑超平面割（
\srcpath{src/optimal\_power\_flow/three\_phase\_hybrid\_relaxation.cpp:violated\_cone\_cuts}，
函数 \srcpath{violated\_cone\_cuts}）：交流提升锥
$\lVert(2W^{R},2W^{I},W_{ii}-W_{jj})\rVert_2\le W_{ii}+W_{jj}$（:append\_if\_violated（lambda））、
PSD 团最小特征向量的秩一割（\srcpath{append\_ac\_psd\_cuts}）、
直流提升锥（:append\_if\_violated（lambda））、换流器视在功率锥（:append\_if\_violated（lambda））、换流器相电流锥
（:append\_if\_violated（lambda））。外逼近循环至多 \fld{max\_outer\_approximation\_rounds}$+1$ 轮，
每轮把累积割并入 LP 交 HiGHS 求解，无违例割时置
\fld{soc\_outer\_approximation\_converged} 并退出（:solve\_three\_phase\_hybrid\_opf\_relaxation；HiGHS 不可用直接
返回 :solve\_three\_phase\_hybrid\_opf\_relaxation）。

\subsubsection*{目标与对偶下界}
LP 目标为 $\sum_g[c_{1,g}S_{\mathrm{base}}p_g + t_g]$，二次成本
$c_2'(p)=c_2S_{\mathrm{base}}^2p^2$ 以 \fld{cost\_tangent\_points} 条支撑切线
$t_g\ge 2c_2'p_0p-c_2'p_0^2$ 在 $[p_{\min},p_{\max}]$ 上均匀布点低估
（\srcpath{src/optimal\_power\_flow/three\_phase\_hybrid\_relaxation.cpp:solve\_three\_phase\_hybrid\_opf\_relaxation}）。
无功正则项与换流器正则项从目标中省略（:solve\_three\_phase\_hybrid\_opf\_relaxation 登记的
\fld{model\_limitations}）。每轮 LP 的对偶按下式构造有效下界：不等式对偶截断到
非正锥后，在变量盒界上逐变量最小化拉格朗日函数（
\srcpath{evaluate\_lp\_dual\_bound}）；各轮取最大写入
\fld{round\_lower\_bounds} 与 \fld{objective\_lower\_bound}（:solve\_three\_phase\_hybrid\_opf\_relaxation）。
由于每条割都是必要条件、成本被低估，该下界对原非凸 OPF 成立
（\fld{outer\_relaxation\_valid} 标记，:solve\_three\_phase\_hybrid\_opf\_relaxation）。

\subsubsection*{输入拒绝与已知省略（诚实口径）}
以下输入被显式拒绝：选项非法（轮数负、切点数 $<2$、容差非正，:solve\_three\_phase\_hybrid\_opf\_relaxation）；
\fld{i\_ac\_fixed} 非零（固定电流注入需要增广电压提升，未实现，:solve\_three\_phase\_hybrid\_opf\_relaxation）；
负二次成本（非凸，:solve\_three\_phase\_hybrid\_opf\_relaxation）。以下模型成分被显式省略并登记：序感知换流器
（跟网/构网）控制等式为非凸，从下界模型中剔除（:solve\_three\_phase\_hybrid\_opf\_relaxation）；目标正则项省略
（:solve\_three\_phase\_hybrid\_opf\_relaxation）；无源性检查失败时聚合网损割不施加（:solve\_three\_phase\_hybrid\_opf\_relaxation）。
因此松弛层输出仅是下界证书与诊断，不给出可行运行点；其
\fld{generator\_active\_power\_pu} 等向量是 LP 解的对应分量，不保证交流可实施。

\subsection{公开选项与结果字段}\label{subsec:tp-fields}

\subsubsection*{问题数据 \fld{ThreePhaseHybridOPFCase}}
\begin{paramtable}
\fld{name} & — & — & — & 空 & 算例名称（日志用）\\
\fld{base\_mva} & $S_{\mathrm{base}}$ & MVA & $>0$ & 1.0 & 标幺基准功率\\
\fld{y\_ac} & $Y_{ac}$ & pu & — & 空 & $n\times n$ 相域复导纳矩阵（稀疏）\\
\fld{i\_ac\_fixed} & $i_{\mathrm{fix}}$ & pu & — & 空 & 固定电流注入；松弛层要求为零\\
\fld{p\_load\_pu} & $p^{\mathrm{load}}$ & pu & — & 空 & 相节点恒功率有功负荷\\
\fld{q\_load\_pu} & $q^{\mathrm{load}}$ & pu & — & 空 & 相节点恒功率无功负荷\\
\fld{v\_min\_pu} & $\underline{v}$ & pu & 0.85--1.0 & 空 & 相节点电压幅值下限\\
\fld{v\_max\_pu} & $\overline{v}$ & pu & 1.0--1.15 & 空 & 相节点电压幅值上限\\
\fld{voltage\_start} & $V^{(0)}$ & pu & — & 空 & 相节点复电压初值\\
\fld{ac\_phase\_index} & — & — & 0/1/2 & 空 & 各相节点所属相号；序感知控制必需\\
\fld{reference\_nodes} & — & — & — & 空 & 交流参考（松弛）相节点列表\\
\fld{reference\_voltage} & $V_{\mathrm{ref}}$ & pu & — & 空 & 参考节点复电压设定\\
\fld{three\_phase\_bus\_nodes} & — & — & — & 空 & 三相齐全母线的节点三元组，VUF 作用集\\
\fld{vuf\_max} & $\varepsilon_{\mathrm{vuf}}$ & pu & $[0,1]$ & 0.03 & 负序/正序电压幅值比上限\\
\fld{generators} & — & — & — & 空 & \fld{PhaseGenerator} 列表（分相机组）\\
\fld{converters} & — & — & — & 空 & \fld{PhaseVSC} 列表\\
\fld{g\_dc} & $G_{dc}$ & pu & — & 空 & 直流电导矩阵（拉普拉斯型）\\
\fld{p\_dc\_load\_pu} & $p^{\mathrm{dc,load}}$ & pu & — & 空 & 直流节点恒功率负荷\\
\fld{v\_dc\_start} & $u^{(0)}$ & pu & — & 空 & 直流电压初值\\
\fld{v\_dc\_min\_pu} & $\underline{u}$ & pu & 0.85--1.0 & 空 & 直流电压下限\\
\fld{v\_dc\_max\_pu} & $\overline{u}$ & pu & 1.0--1.15 & 空 & 直流电压上限\\
\fld{dc\_reference\_terminals} & — & — & — & 空 & 直流电压参考端子\\
\fld{dc\_reference\_voltage\_pu} & $u_{\mathrm{ref}}$ & pu & — & 空 & 直流参考电压设定\\
\end{paramtable}

\subsubsection*{设备结构 \fld{PhaseGenerator} 与 \fld{PhaseVSC}}
\begin{paramtable}
\fld{phase\_node} & — & — & $[0,n)$ & $-1$ & 机组挂接的相节点\\
\fld{p\_min\_pu} / \fld{p\_max\_pu} & $p_{\min}/p_{\max}$ & pu & — & 0 & 有功出力求解界\\
\fld{q\_min\_pu} / \fld{q\_max\_pu} & $q_{\min}/q_{\max}$ & pu & — & 0 & 无功出力求解界\\
\fld{cost\_c2} & $c_2$ & — & $\ge 0$ & 0 & 二次成本系数（按 MW 计；负值被松弛层拒绝）\\
\fld{cost\_c1} & $c_1$ & — & — & 0 & 一次成本系数（按 MW 计）\\
\end{paramtable}

\begin{paramtable}
\fld{phase\_nodes} & — & — & — & 空 & 换流器交流相端子节点（序感知控制须为 A/B/C 三相）\\
\fld{dc\_terminal} & — & — & $[0,n_{dc})$ & $-1$ & 直流端子编号\\
\fld{efficiency} & $\eta$ & pu & $[0.01,1]$（使用时截断） & 0.98 & 交直流功率耦合效率\\
\fld{s\_max\_pu} & $s_{\max}$ & pu & $\ge 0$ & 0 & 三相总视在功率上限\\
\fld{phase\_current\_max\_pu} & $i_{\max}$ & pu & $\ge 0$ & 0 & 相电流上限；非正时回退 $s_{\max}/\sqrt{n_\varphi}$\\
\fld{fixed\_unity\_power\_factor} & — & — & true/false & false & true 时强制 $q^{ac}=0$\\
\fld{control\_mode} & — & — & EqualPhase\allowbreak Power/\allowbreak GridFollowingPLL/\allowbreak GridFormingDroop & EqualPhase\allowbreak Power & 相域控制律；\fld{PositiveSequenceCurrent} 为 \fld{GridFollowingPLL} 别名\\
\fld{nominal\_frequency\_hz} & $f_N$ & Hz & — & 50.0 & 额定频率（当前仅存档，未入方程）\\
\fld{pll\_kp} / \fld{pll\_ki} & $k_p^{\mathrm{pll}}/k_i^{\mathrm{pll}}$ & — & — & 0.01 / 1.0 & PLL 增益；仅用于解后动态平衡点恢复\\
\fld{virtual\_r\_pu} / \fld{virtual\_x\_pu} & $r_v/x_v$ & pu & $>0$（构网型必需） & 0 / 0.10 & 构网诺顿虚拟阻抗\\
\fld{p\_droop\_pu} / \fld{q\_droop\_pu} & — & pu & — & 0.01 / 0.05 & 下垂系数（当前未进入 OPF 方程，仅存档）\\
\fld{voltage\_reference\_pu} & $V_{\mathrm{ref}}^{c}$ & pu & — & 1.0 & 构网电压参考（当前未进入 OPF 方程）\\
\fld{voltage\_integral\_gain} & $k_i^{V}$ & — & — & 10.0 & 电压积分增益；仅用于平衡点恢复\\
\end{paramtable}

\subsubsection*{求解选项 \fld{ThreePhaseHybridOPFOptions}}
\begin{paramtable}
\fld{variant} & — & — & Full/\allowbreak GraphReduced & Full & 全节点或稀疏 Kron 图降阶模型（src/optimal\_power\_flow/three\_phase\_hybrid\_opf.cpp:build\_model\_data）\\
\fld{backend} & — & — & Ipopt/\allowbreak NativeIPM & Ipopt & NLP 后端；Ipopt 受嵌入式编译门控\\
\fld{reduction\_options} & — & — & — & 见下 & 图降阶参数 \fld{SparseKronOptions}：\fld{max\_front}=12、\fld{max\_nnz\_ratio}=2.0、\fld{pivot\_tolerance}=$10^{-12}$、\fld{drop\_tolerance}=$10^{-13}$\\
\fld{max\_iterations} & — & — & $\ge 1$ & 300 & 内点最大迭代数\\
\fld{tolerance} & $\mathrm{tol}$ & — & $>0$ & $10^{-7}$ & 收敛容差；acceptable 级别放宽 10 倍\\
\fld{phase\_one\_time\_limit\_ms} & — & ms & — & 5000 & Phase I 协作式软截止；$\le0$ 表示无限墙钟预算\\
\fld{phase\_one\_max\_iterations} & — & — & $\ge0$ & 12 & Phase I Newton 迭代硬上限\\
\fld{phase\_one\_max\_factorizations} & — & — & $\ge0$ & 14 & primal 恢复与 dual 恢复共用的稀疏数值分解硬上限\\
\fld{phase\_one\_max\_backtracks} & — & — & $\ge0$ & 12 & 每次 Newton 全步之外允许的阻尼回溯数\\
\fld{phase\_one\_barrier\_mu} & $\mu_0$ & — & $(0,0.1]$ & 0.1 & Phase-I central target\\
\fld{phase\_one\_admission\_mu\_factor} & $\eta_a$ & — & $\ge0$ & 1.0 & 局部 Phase I 准入半径 $\eta_a\mu_0$\\
\fld{phase\_one\_primal\_mu\_factor} & $\eta_p$ & — & $\ge0$ & 0.1 & handoff primal corridor $\eta_p\mu_0$\\
\fld{phase\_one\_centrality\_tolerance} & $\eta_c$ & — & $\ge0$ & 0.5 & 最大相对中心乘积偏差\\
\fld{warm\_start\_with\_ipopt} & — & — & true/false & false & 仅无限墙钟预算下启用的旧 Ipopt 可行性热解兼容路径\\
\fld{verify\_derivatives} & — & — & true/false & false & 中心差分校核雅可比/黑塞（src/optimal\_power\_flow/three\_phase\_hybrid\_opf.cpp:solve\_three\_phase\_hybrid\_opf\_impl）\\
\fld{verbose} & — & — & true/false & false & 输出阶段日志与初始化诊断\\
\fld{use\_constraint\_oracle} & — & — & true/false & false & 启用精确约束预言机（src/optimal\_power\_flow/three\_phase\_hybrid\_opf.cpp:solve\_three\_phase\_hybrid\_opf\_impl）\\
\fld{oracle\_initial\_margin} & $m_{\mathrm{init}}$ & — & $\ge 0$ & $10^{-3}$ & 初始筛查裕度\\
\fld{oracle\_activation\_margin} & $m_{\mathrm{act}}$ & — & $\ge 0$ & $10^{-5}$ & 轮次激活裕度\\
\fld{oracle\_probe\_iterations} & — & — & $\ge 0$ & 0 & 预测轮迭代数；0 关闭预测轮\\
\fld{oracle\_max\_rounds} & — & — & $\ge 1$ & 8 & 预言机最大轮数\\
\fld{oracle\_seed\_rows} & — & — & — & 空 & 初始强制不等式行（热启动用）\\
\fld{enforced\_inequality\_rows} & — & — & — & 空 & 高级覆盖：显式强制行集；空=全部不等式（\srcpath{include/hacdcpf/optimal\_power\_flow/three\_phase\_hybrid\_opf.hpp:enforced\_inequality\_rows}）\\
\fld{primal\_start} & $x_0$ & — & — & 空 & 原始初值覆盖（维度匹配且有限才采纳）\\
\fld{equality\_dual\_start} & $\lambda_0$ & — & — & 空 & 等式对偶种子；NaN 行表示待恢复\\
\fld{nonlinear\_inequality\_dual\_start} & $\nu_0$ & — & — & 空 & 非线性不等式对偶种子（全量或强制行两种长度皆可）\\
\fld{nonlinear\_slack\_start} & $s_0$ & — & — & 空 & 非线性不等式松弛种子\\
\end{paramtable}

\subsubsection*{求解结果 \fld{ThreePhaseHybridOPFResult}}
\begin{paramtable}
\fld{converged} & — & — & true/false & false & 求解成功且通过 $10^{-6}$ KKT 审计（src/optimal\_power\_flow/three\_phase\_hybrid\_opf.cpp:solve\_three\_phase\_hybrid\_opf\_impl）\\
\fld{status} / \fld{solver} & — & — & — & 空 & 求解器返回状态串 / 求解器名\\
\fld{iterations} & — & — & — & 0 & 内点迭代数（预言机为累计）\\
\fld{variables} / \fld{equalities} / \fld{inequalities} & — & — & — & 0 & $n_{\mathrm{var}}/n_{\mathrm{eq}}/n_{\mathrm{ineq}}$\\
\fld{enforced\_inequalities} & — & — & — & 0 & 实际强制的非线性不等式行数\\
\fld{constraint\_oracle\_rounds} & — & — & — & 0 & 预言机轮数\\
\fld{constraint\_oracle\_added\_rows} & — & — & — & 0 & 预言机累计补入行数\\
\fld{equality\_jacobian\_nonzeros} / \fld{inequality\_jacobian\_nonzeros} & — & — & — & 0 & 初值处雅可比非零元统计（src/optimal\_power\_flow/three\_phase\_hybrid\_opf.cpp:solve\_three\_phase\_hybrid\_opf\_impl）\\
\fld{max\_omitted\_inequality} & — & — & — & $-\infty$ & 被省略行在解点的最大违反量（src/optimal\_power\_flow/three\_phase\_hybrid\_opf.cpp:solve\_three\_phase\_hybrid\_opf\_impl）\\
\fld{eliminated\_phase\_nodes} & — & — & — & 0 & 图降阶消去的相节点数\\
\fld{objective} & $f$ & — & — & 0 & 目标值（成本币种未定）\\
\fld{runtime\_ms} & — & ms & — & 0 & 求解耗时（序列求解首项含共享建模耗时）\\
\fld{initial\_primal\_residual} / \fld{initial\_dual\_residual} & — & — & — & 0 & 初值原始/对偶残差（对偶初始化后评估）\\
\fld{phase\_one\_initial\_violation} / \fld{phase\_one\_constraint\_violation} & — & pu & — & $+\infty$ & Phase I 前后原坐标最大约束/变量界违反量\\
\fld{phase\_one\_dual\_fit\_residual} & — & — & — & $+\infty$ & 式~\eqref{eq:phase-one-dual-fit} 的 basic-fit 或 SparseQR 正规残差；不是完整 stationarity\\
\fld{phase\_one\_primal\_feasible} / \fld{phase\_one\_in\_handoff\_corridor} & — & — & true/false & false & 最终容差可行 / 有限精度 central handoff 邻域；二者不可混用\\
\fld{phase\_one\_dual\_initialized} & — & — & true/false & false & dual-fit/正性/维度证书\\
\fld{phase\_one\_handoff\_primal\_tolerance} / \fld{phase\_one\_perturbed\_primal\_residual} & — & pu & — & 0 / $+\infty$ & central primal 门槛与实测值\\
\fld{phase\_one\_centrality} / \fld{phase\_one\_barrier\_mu} & — & — & — & $+\infty$ / 0 & 相对中心偏差与目标障碍参数\\
\fld{phase\_one\_budget\_exhausted} / \fld{phase\_one\_termination} & — & — & — & false / not-run & 预算耗尽标记与停止原因\\
\fld{phase\_one\_iterations} / \fld{phase\_one\_factorizations} / \fld{phase\_one\_backtracks} & — & — & — & 0 & Phase I 实际硬预算消耗\\
\fld{phase\_one\_runtime\_ms} / \fld{phase\_one\_linear\_solver} & — & ms / — & — & 0 / unselected & Phase I 耗时与 SparseLU/SparseQR 路径\\
\fld{phase\_two\_start\_requested} / \fld{phase\_two\_start\_accepted} & — & — & true/false & false & Phase II 是否请求并通过独立 central warm-start 审计保留起点\\
\fld{phase\_two\_start\_rejection\_reason} & — & — & — & 空 & central 合同未接受原因\\
\fld{phase\_two\_linear\_solver\_backend} & — & — & — & unselected & Phase II 实际 KKT 后端；零迭代时可保持 unselected\\
\fld{initial\_worst\_equality} / \fld{initial\_worst\_inequality} & — & — & — & $-1$ & 初值最差等式/不等式行号\\
\fld{primal\_residual} / \fld{dual\_residual} / \fld{complementarity} & — & — & — & 0 & 解点 KKT 残差（解点重评估口径）\\
\fld{max\_voltage\_violation} & — & pu & — & 0 & 电压幅值窗口最大越限\\
\fld{max\_vuf} & — & pu & — & 0 & 三相母线最大 $|V_2|/|V_1|$（src/optimal\_power\_flow/three\_phase\_hybrid\_opf.cpp:max\_vuf）\\
\fld{max\_converter\_current\_vuf} & — & pu & — & 0 & 换流器相电流负序/正序比最大值（src/optimal\_power\_flow/three\_phase\_hybrid\_opf.cpp:max\_converter\_current\_vuf）\\
\fld{max\_converter\_current\_loading} & — & pu & — & 0 & 换流器相电流负载率最大值（src/optimal\_power\_flow/three\_phase\_hybrid\_opf.cpp:max\_converter\_current\_loading）\\
\fld{max\_converter\_violation} & — & pu & — & 0 & 换流器容量/电流约束最大越限（src/optimal\_power\_flow/three\_phase\_hybrid\_opf.cpp:solve\_three\_phase\_hybrid\_opf\_impl）\\
\fld{max\_dynamic\_equilibrium\_residual} & — & — & — & 0 & 动态平衡点最大微分残差（src/optimal\_power\_flow/three\_phase\_hybrid\_opf.cpp:recover\_dynamic\_equilibria）\\
\fld{max\_equality\_jacobian\_error} / \fld{max\_inequality\_jacobian\_error} / \fld{max\_lagrangian\_hessian\_error} & — & — & — & 0 & 差分导数校核误差（\fld{verify\_derivatives} 开启时）\\
\fld{primal} & $x$ & — & — & 空 & 原始解向量（模型布局）\\
\fld{equality\_dual} / \fld{inequality\_dual} / \fld{inequality\_slack} & $\lambda/\nu/s$ & — & — & 空 & 对偶与松弛（不等式部分已扩展回全量行）\\
\fld{full\_voltage} & $V_{\mathrm{full}}$ & pu & — & 空 & 恢复后的全节点复电压\\
\fld{dc\_voltage} & $u$ & pu & — & 空 & 直流节点电压\\
\fld{generator\_active\_power\_pu} / \fld{generator\_reactive\_power\_pu} & $p_g/q_g$ & pu & — & 空 & 分相机组出力\\
\fld{converter\_phase\_power\_pu} & $s^{ac}$ & pu & — & 空 & 每台换流器各相复功率\\
\fld{converter\_dc\_power\_pu} & $p^{dc}$ & pu & — & 空 & 每台换流器直流功率\\
\fld{enforced\_inequality\_rows} & — & — & — & 空 & 实际强制不等式行集\\
\fld{constraint\_oracle\_trace} & — & — & — & 空 & 各轮诊断 \fld{ConstraintOracleRoundDiagnostic}（轮次/行数/新增/迭代/收敛/耗时/全量最大违反）\\
\fld{converter\_dynamic\_equilibria} & — & — & — & 空 & 每台换流器的动态平衡点 \fld{PhaseVSCDynamicEquilibrium}（PLL 角/积分器、内电势、功率参考、微分残差）\\
\fld{reduction} & — & — & — & — & 图降阶结果 \fld{SparseKronResult}（含恢复证书）\\
\end{paramtable}

\subsubsection*{松弛层选项与结果}
\begin{paramtable}
\fld{max\_outer\_approximation\_rounds} & — & — & $\ge 0$ & 12 & 外逼近最大轮数（0 表示仅基线 LP）\\
\fld{cost\_tangent\_points} & — & — & $\ge 2$ & 9 & 二次成本支撑切线条数\\
\fld{cone\_tolerance} & — & — & $>0$ & $10^{-6}$ & 锥违例判据\\
\fld{solver\_tolerance} & — & — & $>0$ & $10^{-8}$ & LP 求解容差（保留字段；当前实现未透传给 HiGHS 适配器）\\
\fld{verbose} & — & — & true/false & false & 日志开关（当前实现未使用）\\
\end{paramtable}

\begin{paramtable}
\fld{solved} & — & — & true/false & false & LP 循环正常完成\\
\fld{outer\_relaxation\_valid} & — & — & true/false & false & 外松弛有效性标记（求解完成即置真）\\
\fld{dual\_certificate\_available} & — & — & true/false & false & 对偶下界证书是否可用\\
\fld{soc\_outer\_approximation\_converged} & — & — & true/false & false & 末轮无违例锥割\\
\fld{ac\_passivity\_cut\_applied} / \fld{dc\_passivity\_cut\_applied} & — & — & true/false & false & 聚合网损割是否施加\\
\fld{status} / \fld{solver} & — & — & — & 空 & 状态串 / 求解器名\\
\fld{objective\_lower\_bound} & $\underline{f}$ & — & — & 0 & 原非凸 OPF 目标的有效下界（证书）\\
\fld{lp\_primal\_objective} & — & — & — & 0 & 末轮 LP 原始目标\\
\fld{lp\_primal\_dual\_gap} & — & — & — & 0 & 末轮 LP 原始目标与下界之差\\
\fld{runtime\_ms} & — & ms & — & 0 & 耗时\\
\fld{max\_soc\_violation} & — & — & — & 0 & 末轮最大锥违例\\
\fld{max\_ac\_lift\_violation} / \fld{max\_ac\_psd\_violation} / \fld{max\_dc\_lift\_violation} & — & — & — & 0 & 交流提升/PSD 团/直流提升分项违例\\
\fld{max\_converter\_apparent\_power\_violation} / \fld{max\_converter\_current\_violation} & — & — & — & 0 & 换流器视在功率/相电流锥分项违例\\
\fld{primal\_residual} / \fld{dual\_residual} & — & — & — & 0 & HiGHS 报告的原始/对偶可行度\\
\fld{rounds} / \fld{cuts\_added} & — & — & — & 0 & 实际轮数 / 累计割数\\
\fld{variables} / \fld{equalities} / \fld{inequalities} & — & — & — & 0 & 末轮 LP 规模\\
\fld{primal} & $x$ & — & — & 空 & 末轮 LP 解（提升空间，不保证 AC 可实施）\\
\fld{generator\_active\_power\_pu} / \fld{generator\_reactive\_power\_pu} & $p_g/q_g$ & pu & — & 空 & LP 解中的机组分量\\
\fld{round\_lower\_bounds} & — & — & — & 空 & 各轮下界序列（单调不减）\\
\fld{model\_limitations} & — & — & — & 空 & 省略/近似条目（诚实口径）\\
\end{paramtable}

\subsection{验证与校核}\label{subsec:tp-verify}

模块测试目标为 \fld{test\_three\_phase\_hybrid\_opf}（注册于
\srcpath{tests/CMakeLists.txt:test\_three\_phase\_hybrid\_opf}），全部算例共用一个三母线九节点合成混合系统
（\srcpath{tests/test\_three\_phase\_hybrid\_opf.cpp:make\_small\_hybrid\_case}，
\srcpath{make\_small\_hybrid\_case}）：三母线全相齐全，支路阻抗角一致的串联
导纳 $y=8-\mathrm{j}16$ 与 $y=6-\mathrm{j}12$（:make\_small\_hybrid\_case）；末母线分相负荷
$p=0.08+0.01\varphi$、$q=0.025+0.005\varphi$（:make\_small\_hybrid\_case）；电压窗口
$[0.85,1.10]$~pu（:make\_small\_hybrid\_case）；首母线为平衡参考 $1\angle\{0^\circ,-120^\circ,
+120^\circ\}$（:make\_small\_hybrid\_case）；$\varepsilon_{\mathrm{vuf}}=0.05$（:make\_small\_hybrid\_case）；首母线三相
各挂一台 $c_2=0.1$、$c_1=30$ 机组（:make\_small\_hybrid\_case）；直流两节点、电导 50、节点 1 负荷
0.10、节点 0 为直流参考（:make\_small\_hybrid\_case）；一台 VSC 接末母线三相与直流节点 0，
$\eta=0.98$、$s_{\max}=0.25$（:make\_small\_hybrid\_case）。除注明外各用例后端为 Ipopt、
容差 $10^{-7}$。

\begin{itemize}
  \item \textbf{Full 与图降阶一致性}（:TEST\_CASE("Monolithic phase hybrid OPF Full and GR recover the same solution")）：降阶恰消去 3 个相节点
        （:TEST\_CASE），目标相对偏差不超过 $10^{-6}$（:TEST\_CASE），全电压最大偏差
        小于 $10^{-6}$（:TEST\_CASE），电压越限小于 $10^{-8}$（:TEST\_CASE），
        VUF 不超过 $\varepsilon_{\mathrm{vuf}}+10^{-7}$（:TEST\_CASE），
        换流器电流不平衡度两模型一致到 $10^{-6}$（:TEST\_CASE）。
  \item \textbf{序列求解保留时变负荷节点}（:TEST\_CASE）：第二时段在降阶模型的
        被保留节点上改变负荷，序列解与独立求解目标相对偏差不超过 $10^{-7}$、
        电压偏差小于 $10^{-6}$（:TEST\_CASE）。
  \item \textbf{GFL/GFM 动态平衡点}（:TEST\_CASE）：跟网型正序电流模式下换流器
        电流负序比小于 $10^{-7}$、负载率不超过 $1+10^{-7}$、动态平衡残差小于
        $10^{-7}$（:TEST\_CASE），且解析雅可比/黑塞差分误差小于 $10^{-5}$
        （:TEST\_CASE）；把平衡点交给暂态 \fld{GridFollowingInverter} 后导数最大
        值小于 $10^{-7}$（:TEST\_CASE）。构网型下垂模式同样通过负载率/平衡残差/
        导数校核，内电势幅值大于 0.5~pu，\fld{GridFormingInverter} 导数小于
        $10^{-7}$（:TEST\_CASE）。
  \item \textbf{独立耦合潮流复核}（:TEST\_CASE）：
        \srcpath{make\_three\_phase\_hybrid\_pf\_case} 生成的潮流算例以
        $10^{-9}$ 容差收敛，潮流电压与 OPF 运行点交/直流电压偏差均小于
        $2\times10^{-6}$（:SECTION）。
  \item \textbf{NativeIPM 原-对偶搬运}（:SECTION）：先把图降阶解与对偶按
        保留节点映射搬回 Full 模型（含 NaN 标记的待恢复等式对偶行，
        :TEST\_CASE("Native Full certifies a graph-reduced primal-dual transport")），NativeIPM 从该点出发的初始原始/对偶残差不超过 $10^{-6}$，
        最终 KKT 三项残差均不超过 $10^{-6}$，目标与降阶解相对偏差不超过
        $10^{-6}$（:SECTION）。
  \item \textbf{约束预言机}（:TEST\_CASE）：强制行数严格少于全量且被省略行最大
        违反量低于 $-m_{\mathrm{act}}$（:TEST\_CASE），目标与全行求解相对偏差
        不超过 $10^{-7}$、电压偏差小于 $10^{-6}$、KKT 残差不超过 $10^{-6}$
        （:TEST\_CASE）；以预言机行集与解为种子重解时零新增行且行集不变
        （:TEST\_CASE）。
  \item \textbf{松弛下界证书}（:TEST\_CASE）：零轮基线 LP 的下界不超过非凸解
        目标 $+10^{-7}$（:TEST\_CASE）；10 轮外逼近下界不低于基线、不超过非凸
        解目标 $+10^{-7}$、各轮下界单调（容差 $10^{-9}$）、末轮锥违例不大于
        基线（:TEST\_CASE）。
  \item \textbf{固定电流拒绝}（:TEST\_CASE）：注入非零 \fld{i\_ac\_fixed} 时松弛
        层返回未求解、状态串含 \fld{fixed-current} 且登记
        \fld{model\_limitations}（:TEST\_CASE）。
\end{itemize}

另有两条模块级验证通道：其一，Sioux Falls 交通网 + IEEE 123 节点馈线联合算例
校验 adapter 装配规模与 ID 归因（$\ge 120$ 座三相母线、$\ge 250$ 个相节点、
4 条直流母线、4 台 VSC、$\ge 15$ 台分相机组，\fld{vuf\_max} 选项透传；
\srcpath{tests/test\_sioux\_falls\_ieee123\_case.cpp:TEST\_CASE("Sioux Falls--IEEE 123 case preserves public networks and assumptions")}、:TEST\_CASE("Sioux Falls--IEEE 123 stress controls preserve charging injection")）；其二，
论文复现实验工具 \fld{phase\_hybrid\_opf\_benchmark} 与
\fld{sioux\_falls\_ieee123\_experiment}（后者由 \srcpath{HACDCPF\_HAVE\_OPENDSS}
门控，\srcpath{CMakeLists.txt:phase\_hybrid\_opf\_benchmark、sioux\_falls\_ieee123\_experiment}）提供更大规模的基准与审计，
但不属于默认 ctest 套件。截至当前代码，模块尚无 OpenDSS 等外部工具的 OPF 级
对照算例（潮流层对照不属于本章范围）。

\subsection{边界与局限}\label{subsec:tp-limit}

\begin{itemize}
  \item \textbf{活跃研发}：模块接口与覆盖范围仍在演进；默认后端 Ipopt 受嵌入式
        编译门控，未编译时仅返回 \fld{Unavailable:} 状态（
        \srcpath{MIPSolvers/src/engine/solver/external/adapters.cpp:IpoptAdapter::available}）。
  \item \textbf{网络覆盖}：adapter 不支持 DC/DC 变换器、能量路由器与可控 DC 静态
        发电机（遇之抛错，src/optimal\_power\_flow/three\_phase\_hybrid\_adapter.cpp:build\_three\_phase\_hybrid\_model、391--395）；DC 支路仅电阻型；
        无线路潮流约束、无 OLTC/分接头与并联补偿调节变量；换流器损耗仅单效率
        近似。
  \item \textbf{负荷模型}：仅恒功率节点方程；delta 与电压相关负荷只能以导入
        运行点的恒功率等值进入，已在 \fld{model\_limitations} 声明
        （src/optimal\_power\_flow/three\_phase\_hybrid\_adapter.cpp:build\_three\_phase\_hybrid\_model、474--476）。
  \item \textbf{换流器控制}：序感知模式要求三相端子齐全；构网“下垂”在 OPF
        中仅有诺顿耦合等式，\fld{p\_droop\_pu}/\fld{q\_droop\_pu}/%
        \fld{voltage\_reference\_pu}/\fld{nominal\_frequency\_hz} 不进入任何
        约束（仅存档与平衡点恢复用，src/optimal\_power\_flow/three\_phase\_hybrid\_opf.cpp:build\_model\_data）；松弛层完全省略序感知
        控制等式（src/optimal\_power\_flow/three\_phase\_hybrid\_relaxation.cpp:solve\_three\_phase\_hybrid\_opf\_relaxation）。
  \item \textbf{外部电网}：以 $\pm5\times$ 总负荷的虚拟发电机近似（
        src/optimal\_power\_flow/three\_phase\_hybrid\_adapter.cpp:build\_three\_phase\_hybrid\_model），非无限大母线的阻抗模型。
  \item \textbf{目标函数}：三项正则化权重硬编码（src/optimal\_power\_flow/three\_phase\_hybrid\_opf.cpp:objective）；松弛层
        目标省略正则项且以切线低估二次成本，故其目标值是下界而非近似原目标。
  \item \textbf{初始化}：存在非参考节点机组时交流潮流初始化静默跳过
        （src/optimal\_power\_flow/three\_phase\_hybrid\_opf.cpp:initialize\_ac\_power\_flow、1602--1606），初值质量依赖 \fld{voltage\_start}。
  \item \textbf{松弛层}：拒绝非零固定电流注入与负二次成本
        （src/optimal\_power\_flow/three\_phase\_hybrid\_relaxation.cpp:solve\_three\_phase\_hybrid\_opf\_relaxation）；外逼近轮数耗尽时
        \fld{soc\_outer\_approximation\_converged} 为假，下界仍有效但松弛间隙
        未闭合；LP 解不保证交流可实施。
  \item \textbf{验证覆盖}：完整 OPF/KKT 功能测试仍以三母线合成算例为主；
        另有 360 个三相母线的 Phase I 稀疏规模/预算回归。IEEE 123 规模仅
        adapter 装配校验与实验工具；尚无外部工具 OPF 对照或大规模完整
        Phase II 收敛回归。
\end{itemize}
